% Standard preamble for reslop analysis notes.
% Included via pandoc: --include-in-header=preamble.tex
% Do not edit per-analysis unless you have a specific reason.

% --- Graphics and floats ---
\usepackage{graphicx}
\usepackage{placeins}
\usepackage{float}
\usepackage{needspace}

% Default figure sizing: HEIGHT-based, not width-based.
%
% All figures are produced at figsize=(10,10) — square. But figures
% with colorbars (2D plots, correlation matrices) have extended width
% (e.g., 12x10 effective) due to the colorbar axes. Setting size by
% width makes these figures smaller than 1D plots at the same nominal
% width fraction, because the extra colorbar width eats into the budget.
%
% Setting size by HEIGHT avoids this: a 0.45\linewidth height gives
% consistent visual sizing regardless of colorbars, because the plot
% area height is always 10 inches. Two figures side-by-side at
% height=0.45\linewidth will have matching plot areas even if one has
% a colorbar and the other doesn't.
%
% We unset the default width and set height instead.
\setkeys{Gin}{height=0.45\linewidth,keepaspectratio}

% --- Caption formatting ---
% font=small keeps captions visually distinct from body text.
% Full column width for both figure and table captions.
\usepackage[font=small]{caption}

% --- Tables: prevent page splits ---
% Pandoc generates longtable by default, which splits across pages.
% We strongly discourage this by penalizing breaks inside floats.
% The typesetting agent should also convert longtables to regular
% table floats where the table fits on a single page.
\usepackage{longtable}
\setlength{\LTpre}{0pt}
\setlength{\LTpost}{0pt}

% --- Float placement (relaxed) ---
% LaTeX defaults are too conservative for figure-heavy technical docs.
% These settings allow up to 90% of a page to be floats, preventing
% overfull vbox warnings and figure cutoffs.
\renewcommand{\topfraction}{0.9}
\renewcommand{\bottomfraction}{0.9}
\renewcommand{\textfraction}{0.05}
\renewcommand{\floatpagefraction}{0.6}
\setcounter{topnumber}{3}
\setcounter{bottomnumber}{3}
\setcounter{totalnumber}{6}

% --- Paragraph and page breaking ---
% Prevent widows and orphans (single lines at top/bottom of page).
\widowpenalty=10000
\clubpenalty=10000
% Allow ragged page bottoms rather than stretching vertical glue
% (which creates ugly spacing between paragraphs).
\raggedbottom

% --- Unicode glyph support (tectonic / xelatex / lualatex) ---
%
% The analysis notes deliberately use many non-ASCII glyphs OUTSIDE math
% mode: μ σ ℓ ₄ ⁻¹ ≈ → ⊕ √ ± − × ≤ ≥ and various sub/superscripts. With
% the default Latin Modern fonts these raise "Missing character" warnings
% (xetex/luatex) or hard errors (pdflatex inputenc). tectonic compiles
% with the XeTeX engine, so we load fontspec with a Unicode-rich font
% family. This makes the AN compile without per-run hand-fixing of glyphs.
%
% Guarded with iftex so plain pdflatex still works: under pdflatex we fall
% back to inputenc/textcomp and let newunicodechar (below) map the most
% common glyphs to math-mode equivalents.
\usepackage{iftex}
\ifPDFTeX
  % pdflatex path: byte-level UTF-8 input + Latin Modern + textcomp symbols.
  \usepackage[utf8]{inputenc}
  \usepackage[T1]{fontenc}
  \usepackage{lmodern}
  \usepackage{textcomp}
\else
  % XeTeX / LuaTeX path (tectonic uses XeTeX): fontspec with DejaVu, which
  % covers the Greek, arrows, math-operator and sub/superscript glyphs the
  % notes use. DejaVu Serif/Sans/Mono ship with TeX Live and are bundled
  % with tectonic, so this resolves without a local font install.
  \usepackage{fontspec}
  \setmainfont{DejaVu Serif}
  \setsansfont{DejaVu Sans}
  \setmonofont{DejaVu Sans Mono}
\fi

% Per-glyph fallbacks via newunicodechar. These cover glyphs that DejaVu
% Serif lacks (notably U+2113 SCRIPT SMALL L) and, under pdflatex, map the
% commonly-used Unicode symbols to math-mode equivalents so the document
% compiles regardless of engine. Each definition is robust to both engines.
\usepackage{newunicodechar}

% Script small l (U+2113, "ℓ"): DejaVu Serif has no glyph for this, so map
% it to the math italic \ell on every engine.
\newunicodechar{ℓ}{\ensuremath{\ell}}

% Math operators / relations frequently used in prose and captions.
\newunicodechar{≈}{\ensuremath{\approx}}
\newunicodechar{→}{\ensuremath{\rightarrow}}
\newunicodechar{⊕}{\ensuremath{\oplus}}
\newunicodechar{√}{\ensuremath{\surd}}
\newunicodechar{±}{\ensuremath{\pm}}
\newunicodechar{−}{\ensuremath{-}}          % U+2212 MINUS SIGN
\newunicodechar{×}{\ensuremath{\times}}
\newunicodechar{≤}{\ensuremath{\leq}}
\newunicodechar{≥}{\ensuremath{\geq}}
\newunicodechar{∓}{\ensuremath{\mp}}
\newunicodechar{·}{\ensuremath{\cdot}}      % U+00B7 MIDDLE DOT

% Greek letters seen in prose (μ σ etc.). DejaVu provides these under
% XeTeX, but mapping to math mode guarantees consistent rendering and also
% covers the pdflatex path.
\newunicodechar{μ}{\ensuremath{\mu}}
\newunicodechar{σ}{\ensuremath{\sigma}}
\newunicodechar{α}{\ensuremath{\alpha}}
\newunicodechar{β}{\ensuremath{\beta}}
\newunicodechar{γ}{\ensuremath{\gamma}}
\newunicodechar{δ}{\ensuremath{\delta}}
\newunicodechar{χ}{\ensuremath{\chi}}
\newunicodechar{Δ}{\ensuremath{\Delta}}
\newunicodechar{Σ}{\ensuremath{\Sigma}}

% Subscript / superscript digits and signs used in tokens like m₄ℓ and
% fb⁻¹. Mapped to math sub/superscripts so they render on any engine.
\newunicodechar{₀}{\textsubscript{0}}
\newunicodechar{₁}{\textsubscript{1}}
\newunicodechar{₂}{\textsubscript{2}}
\newunicodechar{₃}{\textsubscript{3}}
\newunicodechar{₄}{\textsubscript{4}}
\newunicodechar{₅}{\textsubscript{5}}
\newunicodechar{₆}{\textsubscript{6}}
\newunicodechar{₇}{\textsubscript{7}}
\newunicodechar{₈}{\textsubscript{8}}
\newunicodechar{₉}{\textsubscript{9}}
\newunicodechar{⁰}{\textsuperscript{0}}
\newunicodechar{¹}{\textsuperscript{1}}
\newunicodechar{²}{\textsuperscript{2}}
\newunicodechar{³}{\textsuperscript{3}}
\newunicodechar{⁴}{\textsuperscript{4}}
\newunicodechar{⁵}{\textsuperscript{5}}
\newunicodechar{⁶}{\textsuperscript{6}}
\newunicodechar{⁷}{\textsuperscript{7}}
\newunicodechar{⁸}{\textsuperscript{8}}
\newunicodechar{⁹}{\textsuperscript{9}}
\newunicodechar{⁻}{\textsuperscript{\ensuremath{-}}}
\newunicodechar{⁺}{\textsuperscript{\ensuremath{+}}}
